\documentclass{beamer}

\title{What is Artificial Intelligence?}

\begin{document}

\begin{frame}
    \titlepage
\end{frame}

\begin{frame}{What is Artificial Intelligence?}
Intelligence \(\rightarrow\) Understand, think, \\
Artificial \(\rightarrow\) Not Natural \& learn
\end{frame}

\begin{frame}{What is Artificial Intelligence?}
\begin{itemize}
\item Intelligence → “the capacity to learn and solve problems” (Websters dictionary)
\item the ability to solve novel problems
\item the ability to act rationally
\item the ability to act like humans
\item Artificial Intelligence → build and understand intelligent entities or agents
\end{itemize}
\end{frame}

\begin{frame}{What is Artificial Intelligence?}
\begin{itemize}
\item “The branch of computer science that is concerned with the automation of intelligent behavior.” (Luger and Stublefield, 1993)
\item AI leverages computers and machines to mimic the problem-solving and decision-making capabilities of the human mind.
\item Systems that think like humans rationally
\item Systems that act like humans rationally
\item In terms of human performance
\item Measure of Success in terms of rationality
\end{itemize}
\end{frame}

\begin{frame}{Acting humanly: The Turing Test approach}
\begin{itemize}
\item "The art of creating machines that perform functions that require intelligence when performed by people"
(Kurzweil, 1990)
\item AI is
\begin{itemize}
\item "The study of how to make computers do things at which, at the moment, people are better."
(Rich and Knight, 1991)
\end{itemize}
\item Turing Test → provide a satisfactory operational definition of intelligence
\item A computer passes the test if a human interrogator, after posing some written questions, cannot tell whether the written responses come from a person or from a computer.
\end{itemize}
\end{frame}

\begin{frame}{Acting humanly: The Turing Test approach}
\begin{itemize}
\item The computer would need to possess the following capabilities to pass the Turing Test:
\begin{itemize}
\item Natural language processing → enables it to communicate successfully
\item Knowledge representation → stores what it knows or hears
\item Automated reasoning → use the stored information to answer questions and to draw new conclusions
\item Machine learning → adapt to new circumstances and detects and extrapolate patterns
\end{itemize}
\end{itemize}
\end{frame}

\begin{frame}{Acting humanly: The Turing Test approach}
\begin{itemize}
\item Total Turing Test → includes a video signal so that the interrogator can test the subject’s perceptual abilities, as well as the opportunity for the interrogator to pass physical objects “through the hatch.”
\item the interrogation should involve visual acuity and physical interaction
\item To pass the total Turing Test, the computer will need:
\begin{itemize}
\item Computer vision → perceive objects
\item Robotics → manipulate objects and move about
\end{itemize}
\end{itemize}
\end{frame}

\begin{frame}{Thinking humanly: Cognitive modeling approach}
\begin{itemize}
\item “The exciting new effort to make computers think . . . machines with minds, in the full and literal sense.” (AI, 1985)
\item “[The automation of] activities that we associate with human thinking, activities such as decision-making, problem solving, learning . . .” (Bellman, 1978)
\item We Must Have Some Way Of Determining How Humans Think
\item Cognitive Science → Interdisciplinary field (AI, psychology, linguistics, philosophy, anthropology) that tries to construct computational theories of human cognition.
\item Try to get “inside” our minds
\end{itemize}
\end{frame}

\begin{frame}{Thinking rationally: Laws of thought approach}
\begin{itemize}
\item “The study of mental faculties through the use of computational models.”
(Charniak and McDermott, 1985)
\item AI is
\begin{itemize}
\item “The study of the computations that make it possible to perceive, reason, and act.”
(Winston, 1992)
\item A system is rational if it does the “right thing,” given what it knows
\begin{itemize}
\item Represent facts about the world via logic
\item Logicist Program → programs could solve any solvable problem described in logical notation.
\item if no solution exists, the program might loop forever
\end{itemize}
\end{itemize}
\end{itemize}
\end{frame}

\begin{frame}{Acting rationally: Rational agent approach}
\begin{itemize}
\item “Computational Intelligence is the study of the design of intelligent agents.”
(Poole et al., 1998)
\item AI is
\begin{itemize}
\item “AI . . . is concerned with intelligent behavior in artifacts.”
(Nilsson, 1998)
\item An agent → an entity that perceives its environment and is able to execute actions to change it, and create and pursue goals.
\item A rational agent → an agent that acts to achieve the best outcome or the best-expected outcome (if there is uncertainty).
\item All the skills needed for the Turing Test also allow an agent to act rationally.
\end{itemize}
\end{itemize}
\end{frame}

\begin{frame}{Foundations of Artificial Intelligence}
\begin{itemize}
\item Many older disciplines contribute to a foundation for artificial intelligence
\item Logic, methods of reasoning, mind as physical system, foundations of learning, language, Philosophy, rationality
\item Formal representation and proof, algorithms, computation, (un)decidability, Mathematics, (in)tractability
\item Modeling uncertainty, learning from data
\item Probability/Statistics, Utility, decision theory, rational economic agents, Economics
\item Neurons as information processing units, Neuroscience
\item How do people behave, perceive, process cognitive information, represent knowledge, Psychology/Cognitive Science
\item Building fast computers, Computer engineering
\item Design systems that maximize an objective function over time, Control theory
\item Knowledge representation, grammars, Linguistics
\end{itemize}
\end{frame}

\begin{frame}{History of Artificial Intelligence}
\begin{itemize}
\item 1940-1950: Early beginnings
\item 1943: McCulloch \& Pitts: Boolean circuit model of brain
\item 1950: Turing's “Computing Machinery and Intelligence”
\item 1950s: Early AI programs, including Samuel's checkers program, Newell \& Simon's Logic Theorist, Gelernter's Geometry Engine
\item 1956: Dartmouth meeting: “Artificial Intelligence” adopted
\item 1965: Robinson's complete algorithm for logical reasoning
\item 1970-1990: Knowledge-based approaches
\item 1969—79: Early development of knowledge-based systems
\item 1980—88: Expert systems industry booms
\item 1988—93: Expert systems industry busts: “AI Winter”
\end{itemize}
\end{frame}

\begin{frame}{History of Artificial Intelligence}
\begin{itemize}
\item 1990—: Statistical approaches
\item Resurgence of probability, focus on uncertainty
\item General increase in technical depth
\item Agents and learning systems… “AI Spring”?
\item 2000—: Where are we now?
\end{itemize}
\end{frame}

\begin{frame}{What can Artificial Intelligence do today?}
\begin{itemize}
\item Robotic vehicles
\item Speech recognition
\item Autonomous planning and scheduling
\item Game playing
\item Spam fighting
\item Logistics planning
\item Robotics
\item Machine Translation
\end{itemize}
\end{frame}

\begin{frame}{Artificial Intelligence in Your Everyday Life}
\begin{itemize}
\item Filters
\item Recommendations
\end{itemize}
\end{frame}

\begin{frame}{Artificial Intelligence in Your Everyday Life}
\begin{itemize}
\item Auto-Complete Sentences
\item Virtual Assistant
\end{itemize}
\end{frame}

\begin{frame}{Artificial Intelligence in Your Everyday Life}
\begin{itemize}
\item Auto-Parking
\end{itemize}
\end{frame}

\end{document}
